\section*{Chapitre \ref{ch:m1:european-equity-market-structure} --- Structure du marché européen des actions}

\begin{solution}{exo:m1:european-equity-market-structure:1}
(a) \omterm{def:m1:european-equity-market-structure:rm}{marché réglementé}. (b) \omterm{def:m1:exchanges-brokers-venues:mtf}{système multilatéral de négociation} (MTF). (c)
\omterm{def:m1:european-equity-market-structure:si}{internalisateur systématique}. (d) \omterm{def:m1:european-equity-market-structure:rm}{système organisé de négociation} (OTF), autorisé parce que les
instruments sont des obligations.
\end{solution}

\begin{solution}{exo:m1:european-equity-market-structure:2}
$H = 0.2025 + 0.0625 + 0.0225 + 0.01 + 0.0025 = 0.30$~; nombre effectif
3.33.
\end{solution}

\begin{solution}{exo:m1:european-equity-market-structure:3}
Enchères de clôture $0.243 \times 57.5 = \text{\euro}14.0$ milliards par jour~;
\omterm{def:m1:european-equity-market-structure:periodic}{enchères périodiques} 6.0~milliards d'euros~; carnets non affichés 6.3~milliards d'euros.
Hors bourse $80 - 57.5 = 22.5$, dont 14.3~milliards d'euros pour les \omterm{def:m1:european-equity-market-structure:si}{internalisateurs
systématiques} et 8.2~milliards d'euros pour les autres.
\end{solution}

\begin{solution}{exo:m1:european-equity-market-structure:4}
Meilleur cours acheteur 18.350 (Z), meilleur cours vendeur 18.355 (Y). Un achat à 18.355 sur Y se fait au
meilleur cours vendeur~: aucun problème cette fois. Si Z avait proposé un cours plus bas, le courtier
n'aurait enfreint aucune \emph{règle de marché} --- il n'existe pas d'interdiction de \omterm{def:m1:us-equity-market-structure:opr}{trade-through} ---
mais il devrait justifier, au titre de sa politique de \omterm{def:m1:european-equity-market-structure:bestex}{meilleure exécution}, pourquoi Z ne figure pas
parmi ses plateformes (coût de connexion face à l'\omterm{def:m1:retail-flow-and-wholesaling:pfof}{amélioration de prix}
attendue). L'obligation porte sur la politique et sur son réexamen, et non sur
chaque exécution.
\end{solution}

\begin{solution}{exo:m1:european-equity-market-structure:5}
$0.07 \times (4\,400 + 400) - 300 = \text{\euro}36$ millions, soit 9~\%
du volume du mois suivant~: une action peut dépasser 7~\% sur un mois tant que l'année
reste en dessous du seuil.
\end{solution}

\begin{solution}{exo:m1:european-equity-market-structure:6}
Avant~: deux ticks de 0.005, soit 0.01~euro, ce qui représente \qty{8.1}{\bp}. Après, la
fourchette minimale vaut un tick, encore 0.01~euro, \qty{8.1}{\bp}~: la
fourchette affichée n'a pas à changer, mais elle vaut désormais \emph{un seul} tick. Les prix qui
se répartissaient sur deux niveaux de chaque côté se concentrent sur un seul~; la file au meilleur prix
devient bien plus longue, la priorité temporelle prend de la valeur, et un \omterm{def:m1:what-a-trading-firm-does:market-maker}{teneur de marché}
gagne un tick plein par aller-retour, mais attend plus longtemps chaque exécution.
\end{solution}

\begin{solution}{exo:m1:european-equity-market-structure:7}
Mois médian du premier dépassement~: 28, et les cinquante trajectoires dépassent le seuil. Avec une
dérive nulle, la part de marché des \omterm{def:m1:us-equity-market-structure:dark}{dark pools} reste proche de 4.5~\%, son niveau de départ, et aucune
trajectoire ne le franchit~: le plafond ne contraint que les tendances, non le bruit, car une somme
sur douze mois lisse le bruit mensuel.
\end{solution}

\begin{solution}{exo:m1:european-equity-market-structure:8}
Il montre que les exécutions ont été légèrement meilleures que le meilleur cours affiché \emph{sur les
trois plateformes que l'algorithme utilise déjà}~: un indice de référence construit à partir de son propre
ensemble d'opportunités, qu'il ne peut guère manquer d'égaler. Il ne montre pas qu'un
meilleur prix n'était pas disponible sur les quatre autres plateformes ou auprès des
\omterm{def:m1:european-equity-market-structure:si}{internalisateurs systématiques}, ce qui pose la question de la \omterm{def:m1:european-equity-market-structure:bestex}{meilleure exécution}. Correctif~: comparer
à un meilleur cours acheteur et un meilleur cours vendeur, construits à partir de toutes les plateformes transparentes
représentant une part significative (ou, désormais, le système consolidé de publication), indiquer la liste des plateformes dans le rapport, et
ajouter le point médian comme second indice de référence, puisqu'une grande part du volume
s'y échange.
\end{solution}

\begin{solution}{pb:m1:european-equity-market-structure:1}
\textbf{1.} 10\,500 millions d'euros par mois, dont 787.5 millions d'euros en \omterm{def:m1:us-equity-market-structure:dark}{dark pool}.
\textbf{2.} $H = 0.36 + 4\times0.01 = 0.40$~: 2.5 opérateurs effectifs.
\textbf{3.} $H = 0.1444 + 0.04 + 0.0729 + 0.0056 + 0.0016 + 0.0012 = 0.266$~:
3.76 mécanismes effectifs.
\textbf{4.} L'usage dépasse 7~\%~: la \omterm{def:m1:european-equity-market-structure:waivers}{dérogation au prix de référence} est suspendue pour
cette action pendant trois mois sur toutes les plateformes.
\textbf{5.} Carnet primaire $38 + 0.2\times7.5\times38/65 = 38.88$~; enchère de
clôture 20~; MTF transparent $27.62$~; \omterm{def:m1:us-equity-market-structure:dark}{dark pool} 0~; \omterm{def:m1:european-equity-market-structure:periodic}{enchères périodiques} $4 + 3.375 = 7.38$~; blocs
$3.5 + 1.875 = 5.38$~; total 99.25 (0.75 non négocié).
\textbf{6.} 39.2~\%, 20.2~\%, 27.8~\%, 0, 7.4~\%, 5.4~\%.
\textbf{7.} $+3.4$ points de pourcentage.
\textbf{8.} 3.57~: moins de mécanismes effectifs, puisque l'un des six a
disparu.
\textbf{9.} Flux en \omterm{def:m1:us-equity-market-structure:dark}{dark pool} de 4.8~millions d'euros par jour~; 30~\% de celui-ci (1.44~million d'euros)
ne reçoit plus le point médian et paie 2 points de base~: 288~euros par jour,
6\,048 euros par mois sur cette seule action.
\textbf{10.} $9 \times 7.5\%/12 = 5.6\%$.
\textbf{11.} Les trois mois vides restent dans la fenêtre jusqu'à ce qu'ils aient
douze mois. L'usage est de 5.6~\% pendant neuf mois, puis 6.25~\%, 6.875~\% et 7.5~\%
au douzième mois suivant le rétablissement~: un cycle de douze mois actifs, trois
mois suspendus.
\textbf{12.} 7~\%.
\textbf{13.} En dessous de 7~\%, personne n'est jamais suspendu et le volume total en \omterm{def:m1:us-equity-market-structure:dark}{dark pool}
sur un cycle est plus élevé ($12\times7 = 84$ contre $12\times7.5 = 90$ suivi
de trois mois sans rien, soit 90 sur quinze mois, ou 72 sur douze mois).
Mais chaque opérateur gagne à prendre davantage de parts tant que la dérogation dure, et
laisserait les autres supporter la contrainte~: un problème de bien commun.
\textbf{14.} Les participants au marché transparent~: le législateur considère que
trop de négociation en \omterm{def:m1:us-equity-market-structure:dark}{dark pool} dégrade le prix public que tout le monde utilise.
\textbf{15.} L'investisseur institutionnel qui perd les exécutions au point médian, et
en fin de compte les épargnants qui se trouvent derrière lui~; une partie du flux d'ordres cesse aussi de négocier.
\textbf{16.} Marché transparent~: un cours indicatif et une quantité sont publiés avant le
croisement, et tout participant peut s'y joindre. \omterm{def:m1:us-equity-market-structure:dark}{Dark pool}~: la publication dure
quelques millisecondes, ne révèle aucun ordre individuel, et le prix est généralement calé
à l'intérieur du cours affiché sur le marché de référence, de sorte que le contenu informationnel est proche
de celui d'un croisement au point médian.
\textbf{17.} Parce qu'afficher un ordre de très grande taille déplacerait le prix
contre lui~; la dérogation protège l'ordre, et les blocs sont trop irréguliers et trop
rares pour se substituer au marché transparent.
\textbf{18.} Plus facile~: tout ce qui se mesure par rapport à un meilleur cours acheteur et vendeur, comme
le gain de fourchette de la question 9 et toute statistique de fourchette effective. Inchangés~:
les parts de marché par plateforme et par mécanisme et le taux d'usage du plafond, qui proviennent des
comptes rendus de transactions, et non des cotations.
\textbf{19.} \textbf{7.4~\%.}
\textbf{20.} Le flux d'ordres est conservé, et il migre vers le substitut légal le plus proche
de ce qu'une règle vient de rendre coûteux.
\end{solution}

\begin{solution}{iq:m1:european-equity-market-structure:1}
Pas de \omterm{def:m1:us-equity-market-structure:opr}{règle de protection des ordres} en Europe~: la \omterm{def:m1:european-equity-market-structure:bestex}{meilleure exécution} est une politique au niveau du courtier,
et non une obligation de routage au niveau de la plateforme. Pas de système consolidé de publication avant 2026, donc pas de
meilleur cours acheteur et vendeur officiel. De nombreuses bourses de cotation et devises, au lieu d'un
marché national unique, la bourse principale dominant ses propres actions.
Des pas de cotation fixés par paliers (prix et liquidité), et non un centime uniforme. Négociation dark
sous dérogation, avec un plafond de volume~; de grandes enchères de clôture et des \omterm{def:m1:european-equity-market-structure:periodic}{enchères périodiques}.
Pas de \omterm{def:m1:retail-flow-and-wholesaling:pfof}{paiement pour flux d'ordres} (PFOF) dans la plupart des juridictions~; négociation bilatérale par
\omterm{def:m1:european-equity-market-structure:si}{internalisateurs systématiques} au lieu des \omterm{def:m1:retail-flow-and-wholesaling:wholesaler}{wholesalers}.
\iqlookfor{le contraste routage et \omterm{def:m1:european-equity-market-structure:bestex}{meilleure exécution} d'abord~; le reste relève du détail.}
\end{solution}

\begin{solution}{iq:m1:european-equity-market-structure:2}
Une entreprise d'investissement qui négocie pour compte propre contre des ordres de clients en dehors d'une
plateforme, de manière organisée, fréquente et substantielle. Chaque transaction a pour
contrepartie l'entreprise, et elle choisit ses clients~; un \omterm{def:m1:exchanges-brokers-venues:mtf}{système multilatéral de
négociation} (MTF) est neutre~: les ordres des participants s'y confrontent selon des
règles non discrétionnaires, et l'opérateur ne prend aucune position.
\iqlookfor{bilatéral et pour compte propre contre multilatéral et neutre.}
\end{solution}

\begin{solution}{iq:m1:european-equity-market-structure:3}
Les fonds indiciels et les gérants indexés sur un indice de référence sont mesurés par rapport à la clôture et veulent
une \omterm{def:m1:the-buy-side:benchmark}{tracking error} nulle par rapport à elle~; les produits dérivés, les fonds et les indices se règlent sur elle~;
et l'enchère offre une grande taille sans fourchette. À mesure que le volume se concentre
là, le carnet d'ordres continu s'amincit, l'impact intrajournalier augmente, et
de plus en plus de participants attendent la clôture~: une migration qui s'auto-alimente et
confère à la bourse de cotation un pouvoir de fixation des prix sur sa propre enchère.
\iqlookfor{la boucle de rétroaction, et pas seulement «~à cause de la \omterm{def:m1:the-buy-side:passive}{gestion passive}~».}
\end{solution}

\begin{solution}{iq:m1:european-equity-market-structure:4}
Les plateformes à inclure (la seule bourse de cotation, toutes les plateformes transparentes, celles qui dépassent un
seuil de part de marché, cotations des \omterm{def:m1:european-equity-market-structure:si}{internalisateurs systématiques} ou non)~; comment
traiter les indications des \omterm{def:m1:european-equity-market-structure:periodic}{enchères périodiques}~; la conversion de devises pour les actions
multicotées~; la synchronisation des horloges entre plateformes situées dans des centres de données différents~; comment
gérer une suspension de cotation ou une enchère sur la bourse principale pendant que les autres plateformes continuent de négocier~;
s'il faut inclure les tailles et quelle taille minimale. Rien de tout cela n'est prescrit,
les choix doivent donc être documentés et maintenus fixes tout au long de l'analyse.
\iqlookfor{l'indice de référence comme choix de modélisation qui doit être déclaré.}
\end{solution}

\begin{solution}{iq:m1:european-equity-market-structure:5}
La liquidité dark au point médian disparaît pendant trois mois. Déplacer la part
passive, économe en fourchette, de l'ordre vers des \omterm{def:m1:european-equity-market-structure:periodic}{enchères périodiques} et des
plateformes de blocs conditionnels sous la \omterm{def:m1:european-equity-market-structure:waivers}{dérogation pour taille élevée}~; prévoir des taux d'exécution plus faibles au
point médian et donc soit un calendrier plus long, soit une exécution plus transparente, qui
paie la fourchette~; recourir davantage à la clôture. Avant la date de suspension, si la prévision est
fiable, accélérer la partie de l'ordre éligible au dark.
\iqlookfor{substitution vers les \omterm{def:m1:european-equity-market-structure:periodic}{enchères périodiques} et les blocs, et le coût en
probabilité d'exécution.}
\end{solution}

\begin{solution}{iq:m1:european-equity-market-structure:6}
Si l'ancienne fourchette valait deux anciens ticks ou plus, elle peut désormais valoir un nouveau tick~: la
fourchette cotée reste identique ou s'élargit jusqu'au nouveau minimum~; la profondeur au meilleur
prix augmente fortement, car les ordres s'empilent sur moins de niveaux~; les temps d'attente dans la file
s'allongent et les annulations diminuent~; l'\omterm{def:m1:retail-flow-and-wholesaling:pfof}{amélioration de prix} à l'intérieur de la fourchette devient
impossible. Pour un \omterm{def:m1:what-a-trading-firm-does:market-maker}{teneur de marché}, le revenu par exécution augmente et le nombre d'exécutions par unité de
temps diminue, et l'avantage passe de la maîtrise de la fixation des prix à la position dans la file et
la rapidité~; la négociation migre vers les plateformes de point médian et les plateformes périodiques qui échappent au
tick. Tester avec une méthode de différence de différences sur le recalcul annuel~:
les actions ayant changé de palier, comparées à des actions semblables qui ne l'ont tout juste pas fait, en comparant
la fourchette, la profondeur, la fourchette réalisée et les parts par plateforme avant et après.
\iqlookfor{la priorité dans la file comme nouvelle ressource rare, et une
stratégie d'identification claire.}
\end{solution}
